\begin{abstract}
Robot foundation models require paired visual, language, proprioceptive, and action supervision at a scale that exposes the physical variability of embodied interaction. Existing robot datasets remain constrained by limited embodiment coverage, heterogeneous storage conventions, weak temporal structure, and evaluation protocols that rarely test transfer outside the training corpus. \textbf{WhiteTiger} is a billion-frame, fully human-teleoperated, multi-platform robot manipulation dataset designed for foundation-model adaptation, cross-embodiment trajectory organization, and event-aware policy learning. WhiteTiger contains \textbf{9521.75 hours} of robot data, \textbf{2989 tasks}, \textbf{513575 episodes}, and \textbf{1028349814 frames} across \textbf{13 physical robot platforms}. The corpus spans fixed-base bimanual manipulators, wheeled mobile manipulators, mobile dual-arm platforms, and humanoid robots, with gripper-based and dexterous-hand end effectors. Platform-specific HDF5 operation records are converted into LeRobot v2.1 while preserving native observation, state, action, and camera schemas.

WhiteTiger couples scalable physical data collection with temporal structure. Human teleoperation grounds each trajectory in real contact dynamics, robot reachability limits, sensor placement, end-effector geometry, and embodiment-specific control spaces. Event-level annotations mark subgoal achievement, phase transitions, boundary crossings, alignment states, handover poses, deformation states, and articulated-object transitions beyond gripper or dexterous-hand actuation. Foundation-model utility is evaluated through offline open-loop zero-shot action prediction using GR00T N1.6 across \textbf{13 platforms} and \textbf{42 paired task-dataset records} from external evaluation datasets excluded from WhiteTiger training. One epoch of full fine-tuning on WhiteTiger reduces \textbf{Normalized Joint MSE by 99.42\%} and \textbf{Normalized ALL MSE by 96.79\%} relative to the no-training checkpoint. The resulting dataset supports studies of robot policy scaling, cross-embodiment adaptation, event-conditioned action generation, hierarchical skill decomposition, progress-aware value learning, and phase-aware evaluation in open-world manipulation settings.
\end{abstract}

\section{Introduction}

Scaling robot policies requires interaction data that couples perception, language, proprioception, and action under the physical constraints of embodied execution. Unlike web-scale image or text corpora, robot trajectories are bounded by robot morphology, actuator limits, camera placement, control frequency, contact dynamics, object affordances, and safety constraints. These factors determine whether policy learning captures transferable sensorimotor structure or memorizes a narrow hardware and scene distribution. Large language models and vision-language models demonstrate predictable benefits from increased data and model scale, but analogous scaling behavior in robotics depends on datasets that preserve the structure of environment-agent interaction rather than flattening embodiment-specific trajectories into weakly aligned logs~\cite{brohan2023rt1,brohan2023rt2,ghosh2024octo,kim2024openvla,black2024pi0}.

Generalist robot policies and vision-language-action models map visual observations, task descriptions, robot states, and prior context to continuous actions~\cite{driess2023palme,brohan2023rt2,kim2024openvla,black2024pi0}. Datasets such as RoboMIND~\cite{wu2024robomind}, DROID~\cite{khazatsky2024droid}, BridgeData V2~\cite{walke2023bridgedatav2}, and Open X-Embodiment~\cite{oneill2023openx} establish the value of real robot trajectories, multi-view observations, task annotations, and standardized interfaces for imitation learning. The remaining bottleneck is not only dataset size. Pretraining-oriented robot corpora must align heterogeneous embodiments, maintain physically meaningful action semantics, expose task diversity across environments and objects, and provide evaluation regimes that test transfer beyond the collected training distribution.

Dense observation-action sequences alone underrepresent the temporal structure of manipulation. Long-horizon and contact-rich tasks contain subgoals, skill boundaries, object-state changes, alignment constraints, and phase transitions that affect downstream policy learning and evaluation. Cloth folding depends on lifting peaks and folding transitions. Insertion depends on pre-insertion alignment and seated states. Container placement depends on boundary crossing. Bimanual handover depends on transfer poses and load exchange. Without explicit temporal structure, these signals must be manually engineered, heuristically recovered, or omitted, which weakens next-event prediction, event-conditioned imitation learning, hierarchical decomposition, skill-boundary detection, future-state prediction, progress-aware reward learning, and phase-aware evaluation.

\textbf{WhiteTiger} targets these constraints through a billion-frame, fully human-teleoperated, multi-platform robot manipulation corpus for embodied foundation model adaptation and event-aware learning. WhiteTiger contains \textbf{9521.75 hours} of data, \textbf{2989 tasks}, \textbf{513575 episodes}, and \textbf{1028349814 frames}. Each trajectory is collected on a physical robot platform and converted from platform-specific HDF5 source records into LeRobot v2.1~\cite{cadene2026lerobot}. The dataset therefore separates scalable training access from native control semantics: LeRobot v2.1 provides a unified loading interface, while platform-specific schemas retain the observation, proprioceptive, action, and camera conventions required for physically grounded learning.

Multi-embodiment coverage is central to the design. WhiteTiger includes fixed-base bimanual manipulators, wheeled mobile manipulators, mobile dual-arm systems, and humanoid robots. The associated control spaces range from arm-and-gripper commands to whole-body action representations involving dexterous hands, head or neck joints, waist motion, legs, and mobile bases. This heterogeneity exposes policies to variation in kinematic topology, reachable workspace, end-effector affordance, visual viewpoint, and actuation granularity. Cross-embodiment trajectory alignment is therefore treated as schema-preserving data organization rather than as an assumption that all robots share an identical action vector~\cite{bousmalis2023robocat,ghosh2024octo,xu2023xskill,liu2024rdt1b}.

Temporal structure is treated as a dataset attribute rather than an auxiliary annotation afterthought. WhiteTiger events identify frames where task progress changes measurably, including subgoal achievement, contact state transitions, object-state changes, alignment states, boundary crossings, handover poses, deformation-related states, and articulated-object transitions. These annotations extend beyond gripper or dexterous-hand opening and closing. The resulting trajectories encode both motor execution and task phase, allowing downstream methods to condition on when subgoals are reached as well as how continuous actions are produced.

The paper presents WhiteTiger through dataset construction, data standardization, and model-oriented evaluation. Dataset construction specifies human teleoperation, platform coverage, task composition, scenario distribution, skill taxonomy, robot configurations, object interactions, and event annotations. Data standardization specifies validation, quality control, metadata normalization, state-action schema alignment, and HDF5-to-LeRobot v2.1 conversion. Model-oriented evaluation uses a GR00T N1.6 checkpoint~\cite{nvidia2025grootn1} under offline open-loop zero-shot action prediction across 13 platforms and 42 paired task-dataset records from datasets external to WhiteTiger and excluded from training. One epoch of full fine-tuning on WhiteTiger reduces Normalized Joint MSE by \textbf{99.42\%} and Normalized ALL MSE by \textbf{96.79\%} relative to the no-training checkpoint.

The contribution of WhiteTiger lies in the joint combination of scale, physical human demonstration, multi-embodiment coverage, schema-preserving standardization, event-level temporal structure, and external model-oriented evaluation. The dataset provides 9521.75 hours, 2989 tasks, 513575 episodes, and 1028349814 frames across 13 physical platforms. Its HDF5-to-LeRobot v2.1 pipeline converts heterogeneous operation records into a training-ready representation while retaining platform-specific action and state semantics. Its morphology, sensing, end-effector, and control-space coverage supports cross-embodiment policy learning from gripper-based manipulation to dexterous and whole-body control. Its event annotations expose task-dependent milestones such as lifting peaks, alignment states, boundary crossings, handover poses, object-state transitions, and deformation states. Its offline zero-shot benchmark quantifies transfer from WhiteTiger fine-tuning to external evaluation datasets rather than measuring only in-corpus reconstruction.

\section{Related Work}

\subsection{Large-scale robot manipulation datasets}

Large-scale robot manipulation datasets provide the empirical substrate for generalist robot policies. Their utility depends on collection setting, embodiment coverage, task diversity, physical interaction diversity, data standardization, and evaluation design. RoboMIND provides 107k demonstration trajectories across 479 tasks and 96 object classes~\cite{wu2024robomind}. DROID collects 76k in-the-wild robot trajectories across 350 hours, 564 scenes, and 84 tasks~\cite{khazatsky2024droid}. BridgeData V2 contains 60,096 trajectories across 24 environments on a low-cost robot platform~\cite{walke2023bridgedatav2}. Open X-Embodiment aggregates data from 22 robots across 21 institutions for cross-embodiment learning at scale~\cite{oneill2023openx}.

WhiteTiger follows this data-centric direction while targeting pretraining-scale physical demonstration, cross-platform schema preservation, event-level temporal annotation, and external foundation-model evaluation. WhiteTiger contains 9521.75 hours of robot data, 2989 tasks, 513575 episodes, and 1028349814 frames. This scale places the dataset in the regime required for studying robot policy scaling under physical interaction constraints rather than only task-specific imitation learning. Recent primitive-level and intermediate-representation datasets further show that temporally structured supervision can support compositional manipulation and plan-then-execute policies~\cite{chen2024rh20tp,li2026robointer,stanovcic2026atlas}.

\begin{table}[t]
\centering
\caption{Comparison with representative large-scale robot manipulation datasets. ``Event ann.'' denotes dataset-level event or keyframe temporal milestone annotations. ``n/r'' indicates that the corresponding statistic is not reported as a primary dataset-level quantity in the cited paper. Task and skill definitions follow the conventions of each source and are therefore not directly interchangeable across datasets.}
\label{tab:dataset-comparison}
\footnotesize
\setlength{\tabcolsep}{4pt}
\renewcommand{\arraystretch}{1.08}
\begin{adjustbox}{max width=\textwidth}
\begin{tabular}{lrrrrrcl}
\toprule
Dataset & Traj. / Ep. & Hours & Platforms & Tasks & Reported skills & Event ann. & Collection \\
\midrule
BridgeData V2~\cite{walke2023bridgedatav2} & 60.1k & n/r & 1 & n/r & 13 & No & Human + scripted \\
DROID~\cite{khazatsky2024droid} & 76k & 350 & 1 & 84 & 86 & No & Human teleoperation \\
RoboMIND~\cite{wu2024robomind} & 107k & n/r & 4 & 479 & 38 & No & Human teleoperation \\
Open X-Embodiment~\cite{oneill2023openx} & $>1$M & n/r & 22 & 160266 & 527 & No & Aggregated datasets \\
\textbf{WhiteTiger} & \textbf{513.6k} & \textbf{9521.75} & \textbf{13} & \textbf{2989} & \textbf{64} & \textbf{Yes} & \textbf{Fully human teleoperation} \\
\bottomrule
\end{tabular}
\end{adjustbox}
\end{table}

Table~\ref{tab:dataset-comparison} identifies the combined position of WhiteTiger. Relative to single-platform corpora such as BridgeData V2 and DROID, WhiteTiger expands embodiment coverage and platform heterogeneity. Relative to RoboMIND, WhiteTiger contains more trajectories, tasks, and physical robot platforms. The 64-skill taxonomy reports behavioral coverage beyond task-count statistics. Open X-Embodiment covers more robot embodiments and a larger aggregate trajectory corpus through institutional aggregation. WhiteTiger instead provides a coherently managed corpus of fully human-teleoperated physical demonstrations across 13 platforms, standardized into LeRobot v2.1, with reported event-level temporal milestones and external model-oriented evaluation.

\subsection{Event-aware robot learning}

Event-aware robot learning uses sparse temporal structure to reduce effective horizon length and expose task progress. PerAct, RVT, and Act3D predict 3D keyframe or keypose actions for manipulation, while waypoint-based and hybrid-action imitation learning methods use sparse intermediate targets to anchor continuous control~\cite{shridhar2022peract,goyal2023rvt,gervet2023act3d,shi2023waypointbased,belkhale2023hydra}. Hierarchical and diffusion-based methods use sparse poses, waypoints, or skills as high-level trajectory anchors with lower-level policies generating motion between them~\cite{ke2024diffuseractor,ma2024hdp,liang2023skilldiffuser}. Skill segmentation and temporal abstraction methods decompose long-horizon demonstrations into subtasks, options, or primitive skills, making event boundaries natural supervision for hierarchical policies and diagnostic evaluation~\cite{shiarlis2018taco,pertsch2020spirl,mao2024dexskills}.

Intermediate-representation methods expose additional structure between task instructions and low-level actions. Vision-language-action systems introduce language motions, embodied reasoning traces, future visual states, contact phases, or per-frame intermediate labels~\cite{belkhale2024rth,zawalski2024ecot,zhao2025cotvla,li2025switchvla,li2026robointer}. Video-planning and generative robot models use visual subgoals or future frames as planning targets~\cite{du2023unipi,wu2023gr1,zhao2025cotvla}. Reward and value learning methods use visual or language-aligned progress signals to shape policy learning~\cite{zakka2021xirl,ma2022vip,ma2023liv,sontakke2023roboclip}. WhiteTiger supplies event annotations as general temporal structure supervision for next-event prediction, event-to-event action generation, hierarchical policy learning, skill-boundary detection, VLA intermediate supervision, future-state prediction, progress-aware reward learning, and phase-aware evaluation.

\subsection{Multi-embodiment robot learning}

Multi-embodiment learning trains policies that benefit from trajectories collected on heterogeneous robot platforms. The central difficulty is not platform count alone. Robot datasets differ in kinematic topology, joint dimensionality, end-effector design, gripper conventions, dexterous-hand actuation, camera placement, control frequency, coordinate frames, and task distribution. Prior work has studied multi-robot datasets, cross-embodiment policy learning, and adaptation to unseen robot configurations~\cite{dasari2020robonet,oneill2023openx,bousmalis2023robocat,ghosh2024octo,xu2023xskill,liu2024rdt1b}.

WhiteTiger provides multi-embodiment trajectories across fixed-base manipulation, wheeled mobile manipulation, mobile dual-arm operation, and humanoid whole-body control. Its platforms use different sensing arrangements, parallel grippers, dexterous hands, and heterogeneous action-state schemas. The dataset is therefore suited to studying schema-preserving robot data standardization, cross-platform policy training, and action prediction under embodiment-dependent physical constraints.

\subsection{Vision-language-action models}

Vision-language-action models learn sensorimotor mappings from visual observations, language instructions, robot state, and context to robot actions. Recent embodied models combine multimodal foundation models with continuous robot control representations to support generalist policies~\cite{driess2023palme,brohan2023rt2,ghosh2024octo,kim2024openvla,black2024pi0,nvidia2025grootn1}. Their generalization depends on synchronized observation-action trajectories, reliable task-level annotations, physically meaningful action labels, and sufficient variation across environments, objects, and embodiments.

WhiteTiger supports this paradigm by organizing human-teleoperated robot trajectories as standardized observation-action sequences in LeRobot v2.1~\cite{cadene2026lerobot}. The benchmark evaluates this data foundation through offline open-loop zero-shot action prediction using GR00T N1.6 checkpoints~\cite{nvidia2025grootn1} on datasets external to WhiteTiger.

\subsection{Data standardization for robot learning}

Robot data standardization is a prerequisite for large-scale policy learning across heterogeneous platforms. Raw operation logs often differ in state definitions, action dimensions, camera streams, gripper conventions, control rates, timestamp conventions, and metadata schemas. A unified policy cannot be trained reliably when state-action dimensions are misaligned, camera streams are ambiguously named, or task labels are inconsistent~\cite{dasari2020robonet,oneill2023openx,ghosh2024octo}.

WhiteTiger addresses this problem through structured dataset management and HDF5-to-LeRobot v2.1 conversion. The management layer tracks data additions, cleaning, normalization changes, and embodiment-specific schema fixes. The LeRobot v2.1 representation provides a training-oriented interface for scalable loading, sharing, and evaluation. This design is critical in pretraining-scale robot corpora because small state-action mapping errors can propagate across many training runs and distort measured scaling behavior.

\subsection{Positioning of WhiteTiger}

WhiteTiger is positioned as a billion-frame, fully human-teleoperated, multi-platform robot manipulation dataset for embodied foundation model adaptation and event-aware robot learning. Its defining property is the joint presence of physical demonstration scale, multi-embodiment morphology, platform-specific schema preservation, HDF5-to-LeRobot v2.1 standardization, event-level temporal annotation, and external offline evaluation with a foundation-model checkpoint. The dataset complements single-hardware corpora and aggregated multi-institution datasets by providing a coherently organized, temporally structured, and model-oriented data foundation for open-world manipulation policy learning.

\section{WhiteTiger Dataset}

\subsection{Dataset scale}

WhiteTiger contains \textbf{9521.75 hours} of robot manipulation data, \textbf{2989 tasks}, \textbf{513575 episodes}, and \textbf{1028349814 frames} across \textbf{13 robot platforms}. The corpus targets foundation-model adaptation and policy scaling rather than isolated task-specific imitation.

The total frame count and duration imply an average recording rate of approximately 30 frames per second. The average episode contains approximately 2002 frames, corresponding to approximately 66.7 seconds. These quantities define the temporal density available for learning continuous sensorimotor mappings and the episode-level horizon distribution that constrains open-loop and closed-loop policy evaluation.
